% !PRINTED 12
\weberpage{12}

We mention further the following theorem, whose correctness follows immediately from
the isomorphism of transformed groups:

\gesperrt{If $Q$ is a divisor of $P$, and $R$ a divisor of $Q$, then, if $Q'$
and $R'$ arise from $Q$ and $R$ by transformation by the same element of $P$,
$R'$ is also a divisor of $Q'$; and if $R$ is a normal divisor of $Q$,
then $R'$ is also a normal divisor of $Q'$.}

\begin{center}
\S.\,4.

\gesperrt{Composition of the Parts.}
\end{center}

If $A$ and $B$ are any two sequences of elements from a group $P$
(groups or not), then by the symbolic product $AB$ we shall understand
the totality of all elements that one obtains when each element $a$ of $A$
is compounded with each element $b$ of $B$, according to the rule holding
in $P$, into $ab$. This manner of composition of
$A$ and $B$ into $AB$ we shall call the \gesperrt{composition of the parts} (of $P$).
We distinguish here between a \gesperrt{part} and a \gesperrt{divisor} of
$P$, so that a divisor is always to be a group, which is not necessary
for a part.

One can likewise compound three and more parts $A, B, C \ldots$ of $P$,
and the associative law holds for the composition of the parts, as
follows immediately from the fact that this law holds in $P$. Accordingly the
meaning of a compositum $ABC \ldots$ is uniquely determined.

In the composition $AB$ one of the parts $A$, $B$ may also consist of a
single element, and then the notation $AB$ agrees with the notation
used above for the subsidiary groups.

For the composition of the parts the following theorems hold:

\noindent 1.\ \gesperrt{The necessary and sufficient condition
for a part $A$ of $P$ to be a group consists in the
equation}
\weq{1}{AA = A.}
For this equation says nothing else than that, if $a$ and $a'$ are contained in
$A$, then $aa'$ also belongs to $A$, and hence that $A$ is a group.
